%=============================================================================
\section{Discussion and Outlook}
\label{sec:discussion}

\subsection{Relation to previous work}
\label{sec:previous}

The connection between nearest-neighbor distances and the void probability function was recognized by \citet{Kerscher1999}.  \citet{Banerjee2020} developed the full $k$NN-CDF framework and demonstrated its sensitivity to all connected $N$-point functions.  \citet{Banerjee2021} extended it to cross-correlations between two point sets.  The present work differs in focus: rather than using $k$NN-CDFs as a beyond-two-point summary statistic for cosmological parameter estimation, we use them as a \emph{computational shortcut} for the standard two-point correlation function measurement.  The goal is not to extract additional information (though the $k$NN-CDFs do contain higher-order information), but to extract the \emph{same} two-point information faster.

This distinction is important for the intended audience.  Survey analysis pipelines are built around $\xi(r)$ and $P(k)$; they have well-characterized covariance models, window function corrections, and systematic error budgets.  Introducing a new summary statistic requires re-deriving all of these.  By showing that the $k$NN approach gives the same $\xi(r)$, we allow it to be adopted as a drop-in replacement for the pair-counting step, without changing anything downstream.

\subsection{Angular correlation functions}
\label{sec:angular}

The formalism developed here applies directly to 3D correlation functions measured from spectroscopic surveys.  For photometric surveys, the primary observable is the angular correlation function $w(\theta)$, measured from projected positions on the sky.  The $k$NN framework generalizes straightforwardly to 2D: the tree is built on angular coordinates, the CDFs measure angular nearest-neighbor distances, and the estimator for $w(\theta)$ follows the same DD/DR ratio structure.  The computational savings are even larger in 2D, since the pair count grows as $N \cdot N_{\rm annulus}$ with $N_{\rm annulus} \propto \theta$ rather than $r^2$.

\subsection{Covariance estimation}
\label{sec:covariance}

Accurate covariance matrices for $\xi(r)$ measurements are critical for cosmological parameter inference.  The current state of the art for DESI uses the RascalC code \citep{Philcox2020,Rashkovetskyi2024}, which computes a semi-analytical covariance in the Gaussian approximation (using the measured $\xi$ and the random catalog to encode the survey geometry) and approximates non-Gaussian contributions by rescaling the shot-noise terms with a single parameter $\alpha_{\mathrm{SN}}$, calibrated from jackknife resampling or a small set of mocks.

The $k$NN framework interacts with this covariance pipeline at three levels:

\paragraph{(i) Drop-in compatibility.}  Since the $k$NN estimator gives the same $\xi(r)$ as LS (Section~\ref{sec:equivalence}), any existing covariance matrix is directly applicable.  No changes to the downstream pipeline are needed.

\paragraph{(ii) Dilution-ladder variance.}  The multiple subsamples at each dilution level provide an empirical covariance estimate, analogous to jackknife but with natural scale separation (Section~\ref{sec:variance}).  This can serve as an independent cross-check on mock-based or semi-analytical covariances.

\paragraph{(iii) Scale-dependent $\alpha_{\mathrm{SN}}$ from kNN residuals.}  As developed in Section~\ref{sec:nonpoisson}, the non-Poissonian residuals $\Delta_k(r_0)$ provide a direct measurement of the excess variance $\delta\sigma^2_N(V)$ at each scale, replacing the single-parameter $\alpha_{\mathrm{SN}}$ with a data-driven, scale-dependent function.  This is the most novel contribution of the kNN approach to the covariance problem: it upgrades the least constrained ingredient of the RascalC pipeline at no additional computational cost.

\subsection{Computational cost}
\label{sec:cost}

A fair comparison with current survey pipelines must account for the optimizations already in use.  The Euclid 2PCF pipeline \citep{EuclidLXXII} uses tree-accelerated pair counting (kd-tree, octree, or linked-list methods) with OpenMP parallelization, split random catalogs to reduce the RR cost, and the Linear Construction method \citep{Keihanen2022} to speed up mock-based covariance estimation.  We compare the asymptotic scaling, but caution that wall-clock performance depends on constants, cache behavior, and parallelization efficiency that can only be determined by benchmarking (Section~\ref{sec:validation}).

\begin{center}
\renewcommand{\arraystretch}{1.15}
\begin{tabular}{@{}lcc@{}}
\toprule
Operation & Tree-accelerated LS & $k$NN ladder \\
\midrule
DD & $O(N_d \cdot N_{\rm shell})$ & $O(N_d \log N_d)$ per level \\
DR & $O(N_R \cdot N_{\rm shell})$ & $O(N_R \log N_d)$ per level \\
RR (uniform, unweighted) & $O(N_R^2)$ or split & Analytic (Erlang) \\
RR (weighted survey) & $O(N_R^2)$ or split & $O(N_R \log N_R)$ \\
Mock-based covariance & $N_{\rm mock} \times$ above & $N_{\rm mock} \times$ above \\
Binning & Required (pre-chosen) & Not needed \\
Non-Poissonian diagnostics & Not available & Free (from same query) \\
\bottomrule
\end{tabular}
\end{center}

Here $N_{\rm shell}$ is the mean number of objects in a radial shell at the separation of interest.  For tree-accelerated pair counting, the effective $N_{\rm shell}$ depends on the tree pruning efficiency and is typically much smaller than the naive shell count, particularly at small separations where the tree prunes aggressively.  At large separations ($r \gtrsim 100\,h^{-1}\,\mathrm{Mpc}$), $N_{\rm shell}$ grows as $r^2$ and tree acceleration provides diminishing returns; this is exactly where the kNN dilution ladder is most advantageous.

The kNN approach does not eliminate the need for mock-based covariance estimation.  Mock catalogs are essential for validating the pipeline, characterizing systematic errors, and computing the model-dependent covariance at different points in parameter space.  Whatever speedup the kNN method provides per catalog applies equally to the mock runs.  The dilution-ladder variance (Section~\ref{sec:variance}) is a complementary, data-driven diagnostic that provides an independent cross-check on the mock-based covariance but does not replace it.

Definitive wall-clock benchmarks at survey scale ($N_d \sim 2 \times 10^7$, $N_R \sim 10^9$) comparing bosque-based kNN queries to Corrfunc and the Euclid 2PCF code are deferred to the companion code paper.

\subsection{Summary}
\label{sec:summary_final}

We have shown that the $k$-Nearest Neighbor CDF framework provides an estimator for the two-point correlation function $\xi(r)$ that is:
\begin{itemize}[nosep]
  \item Mathematically equivalent to the Landy--Szalay estimator when the summed pair-count density is used with sufficient $k_{\max}$.
  \item Computationally faster by orders of magnitude, scaling as $O(N_R \log N_d)$ rather than $O(N^2)$.
  \item Naturally continuous in $r$, requiring no choice of radial binning.
  \item Compatible with survey geometry and per-object weights through the value-weighted $k$NN-CDF formalism.
  \item Extendable to large scales via the dilution ladder, with built-in overlap consistency checks and empirical variance estimation.
  \item Equivalent in statistical precision to pair counting, with the same cosmic variance floor and the same shot noise.
\end{itemize}

Beyond estimating $\xi(r)$ itself, the $k$NN framework provides---at no additional computational cost---a scale-dependent diagnostic of non-Poissonian structure in the galaxy point process (Section~\ref{sec:nonpoisson}).  The residual between the measured kNN-CDFs and the Poisson$|\xi$ prediction yields the excess factorial moments at each scale, directly measuring the non-Gaussian and non-Poissonian contributions that enter the covariance matrix.  This enables a data-driven, scale-dependent replacement for the single-parameter shot-noise rescaling used in current semi-analytical covariance estimation (RascalC), potentially improving covariance accuracy for both BAO and full-shape analyses.

The method is immediately applicable to any survey that currently uses pair counting for $\xi(r)$ estimation.  It requires only a standard KD-tree implementation and the survey random catalog.  We encourage the LSS community to validate the approach on existing mock catalogs and adopt it where the computational savings are significant.

Code implementing the $k$NN ladder estimator, including survey weight support, the dilution scheme, and the non-Poissonian diagnostic, will be made publicly available.
